%% ============================================================
%%  DOCUMENT DE CADRAGE LOGIQUE — AGENT HYPOTHÈSES
%%  3 Chapitres : Introduction & Contexte | Notions Fondamentales | Veille Technologique
%% ============================================================
\documentclass[12pt, a4paper]{report}

%% ---------- Encodage & Langue ----------
\usepackage[utf8]{inputenc}
\usepackage[T1]{fontenc}
\usepackage[french]{babel}

%% ---------- Police propre et professionnelle ----------

%% ---------- Mise en page ----------
\usepackage[top=2.5cm, bottom=2.5cm, left=2.8cm, right=2.8cm]{geometry}
\usepackage{parskip}
\setlength{\parindent}{0pt}
\setlength{\parskip}{6pt}

%% ---------- Couleurs ----------
\usepackage[dvipsnames,svgnames,table]{xcolor}
\definecolor{navyblue}{HTML}{1A3A5C}
\definecolor{midblue}{HTML}{2E6DA4}
\definecolor{lightblue}{HTML}{D6E8F7}
\definecolor{paleblue}{HTML}{EBF4FB}
\definecolor{darkgreen}{HTML}{1B5E3B}
\definecolor{midgreen}{HTML}{2E7D52}
\definecolor{lightgreen}{HTML}{D4EDDA}
\definecolor{codebg}{HTML}{F3F6FA}
\definecolor{stepbg}{HTML}{EBF4FB}
\definecolor{rulegray}{HTML}{B0BEC5}

%% ---------- Titres personnalisés ----------
\usepackage{titlesec}
\titleformat{\chapter}[display]
  {\normalfont\huge\bfseries\color{navyblue}}
  {\color{midblue}\chaptertitlename\ \thechapter}
  {10pt}
  {\color{navyblue}}
  [\vspace{4pt}\textcolor{midblue}{\hrule height 1.5pt}]

\titleformat{\section}
  {\normalfont\Large\bfseries\color{midblue}}
  {\thesection}{0.8em}{}
  [\vspace{2pt}\textcolor{lightblue}{\hrule height 0.8pt}]

\titleformat{\subsection}
  {\normalfont\large\bfseries\color{darkgreen}}
  {\thesubsection}{0.7em}{}

\titleformat{\subsubsection}
  {\normalfont\normalsize\bfseries\color{midgreen}}
  {\thesubsubsection}{0.6em}{}

%% ---------- En-têtes / Pieds de page ----------
\usepackage{fancyhdr}
\pagestyle{fancy}
\fancyhf{}
\fancyhead[L]{\small\textcolor{midblue}{Agent Hypothèses — Document de Cadrage}}
\fancyhead[R]{\small\textcolor{midblue}{\leftmark}}
\fancyfoot[C]{\textcolor{navyblue}{\thepage}}
\renewcommand{\headrulewidth}{0.4pt}
\renewcommand{\headrule}{\hbox to\headwidth{\color{midblue}\leaders\hrule height \headrulewidth\hfill}}

%% ---------- Tableaux ----------
\usepackage{booktabs}
\usepackage{tabularx}
\usepackage{longtable}
\usepackage{array}
\usepackage{multirow}
\usepackage{colortbl}

\newcolumntype{L}[1]{>{\raggedright\arraybackslash}p{#1}}
\newcolumntype{C}[1]{>{\centering\arraybackslash}p{#1}}

%% ---------- Boîtes colorées ----------
\usepackage{tcolorbox}
\tcbuselibrary{skins, breakable, listingsutf8}

%% Boîte formule / règle
\newtcolorbox{formulebox}[1][]{
  colback=paleblue, colframe=midblue, 
  fonttitle=\bfseries\color{navyblue},
  boxrule=0.8pt, arc=4pt,
  left=8pt, right=8pt, top=6pt, bottom=6pt,
  #1
}

%% Boîte étape (fichier expliqué)
\newtcolorbox{stepbox}[2][]{
  colback=stepbg, colframe=navyblue,
  fonttitle=\bfseries\color{white},
  coltitle=white, attach boxed title to top left={yshift=-2mm, xshift=6mm},
  boxed title style={colback=navyblue, colframe=navyblue, arc=3pt},
  title={#2},
  boxrule=0.6pt, arc=4pt,
  left=8pt, right=8pt, top=14pt, bottom=6pt,
  breakable,
  #1
}

%% Boîte résultat / note verte
\newtcolorbox{resultbox}[1][]{
  colback=lightgreen, colframe=darkgreen,
  boxrule=0.6pt, arc=4pt,
  left=8pt, right=8pt, top=6pt, bottom=6pt,
  #1
}

%% Boîte code
\newtcolorbox{codebox}[1][]{
  colback=codebg, colframe=rulegray,
  boxrule=0.5pt, arc=3pt,
  left=6pt, right=6pt, top=4pt, bottom=4pt,
  fontupper=\small\ttfamily,
  #1
}

%% ---------- Listes ----------
\usepackage{enumitem}
\setlist[itemize,1]{label=\textcolor{midblue}{\textbullet}, leftmargin=1.5em, itemsep=2pt}
\setlist[itemize,2]{label=\textcolor{darkgreen}{$\circ$}, leftmargin=1.5em, itemsep=1pt}
\setlist[enumerate,1]{label=\textcolor{navyblue}{\arabic*.}, leftmargin=1.5em, itemsep=3pt}

%% ---------- Liens ----------
\usepackage[hidelinks]{hyperref}

%% ---------- Divers ----------
\usepackage{graphicx}
\usepackage{amsmath}
\usepackage{amssymb}
\usepackage{setspace}
\usepackage{float}

%% ---------- Table des matières ----------
\usepackage{tocloft}
\renewcommand{\cfttoctitlefont}{\Large\bfseries\color{navyblue}}
\renewcommand{\cftchapfont}{\bfseries\color{navyblue}}
\renewcommand{\cftsecfont}{\color{midblue}}
\renewcommand{\cftsubsecfont}{\color{darkgreen}}
\setlength{\cftbeforechapskip}{6pt}

%% ---------- Arbre de fichiers (verbatim coloré) ----------
\usepackage{fancyvrb}

%% ============================================================
\begin{document}
\onehalfspacing

%% ============================================================
%%  PAGE DE GARDE
%% ============================================================
\begin{titlepage}
\pagecolor{navyblue}
\color{white}
\begin{center}

\vspace*{1.5cm}

{\small\bfseries\color{lightblue} UNIVERSITÉ ABDELMALEK ESSAÂDI — FSTT}\\[4pt]
{\small\color{lightblue} Licence Data Analytics}\\[6pt]
{\small\color{lightblue} Projet de Fin d'Études}

\vspace{1.2cm}
\textcolor{midblue}{\hrule height 1pt}
\vspace{0.3cm}
\textcolor{rulegray}{\hrule height 0.5pt}
\vspace{1.2cm}

{\fontsize{13}{18}\selectfont\bfseries\color{lightblue}
DOCUMENT DE CADRAGE LOGIQUE}\\[10pt]

{\fontsize{26}{34}\selectfont\bfseries\color{white}
Agent Hypothèses}\\[14pt]

{\fontsize{14}{20}\selectfont\color{lightblue}
Module Pédagogique Intelligent}\\[6pt]

{\fontsize{11}{16}\selectfont\color{rulegray}
Plateforme IA End-to-End :
Idéation $\rightarrow$ Business Plan $\rightarrow$ Immatriculation}

\vspace{1.2cm}
\textcolor{rulegray}{\hrule height 0.5pt}
\vspace{0.3cm}
\textcolor{midblue}{\hrule height 1pt}
\vspace{1.8cm}

%% Trois blocs thématiques
\begin{minipage}{0.3\textwidth}
\centering
\textcolor{midblue}{\large\bfseries Chapitre 1}\\[4pt]
\textcolor{lightblue}{\small Introduction \&\\Contexte Stratégique}
\end{minipage}
\hfill
\begin{minipage}{0.3\textwidth}
\centering
\textcolor{midblue}{\large\bfseries Chapitre 2}\\[4pt]
\textcolor{lightblue}{\small Notions Fondamentales\\\& Logique Métier}
\end{minipage}
\hfill
\begin{minipage}{0.3\textwidth}
\centering
\textcolor{midblue}{\large\bfseries Chapitre 3}\\[4pt]
\textcolor{lightblue}{\small Veille Technologique\\\& Stack Technique}
\end{minipage}

\vfill

{\small\color{rulegray} Année Universitaire 2025--2026}

\end{center}
\end{titlepage}

\nopagecolor
\color{black}

%% ============================================================
%%  TABLE DES MATIÈRES
%% ============================================================
\tableofcontents
\newpage

%% ============================================================
%%  CHAPITRE 1 — INTRODUCTION & CONTEXTE STRATÉGIQUE
%% ============================================================
\chapter{Introduction \& Contexte Stratégique}

\section{Présentation générale}

Le présent document constitue le cadrage logique du \textbf{module Agent Hypothèses}, composante centrale de la plateforme intelligente d'accompagnement à la création d'entreprise au Maroc, développée dans le cadre d'un Projet de Fin d'Études (PFE) de Licence Data Analytics à la FSTT.

La plateforme repose sur une architecture multi-agents orchestrée par \textbf{n8n}, où chaque agent couvre une phase spécifique du parcours entrepreneurial : idéation, analyse de marché, collecte d'hypothèses, génération financière, conformité juridique et production du Business Plan final.

L'\textbf{Agent Hypothèses}, objet de ce document, constitue le pont stratégique entre le dialogue avec l'entrepreneur et la génération des tableaux financiers prévisionnels.

\section{Objectif du document}

Ce document présente l'architecture technique, les notions fondamentales, la démarche méthodologique et la veille technologique de l'Agent Hypothèses. Il couvre :

\begin{itemize}
  \item La démarche Data Analytics suivie pour concevoir et réaliser l'agent
  \item Les notions fondamentales utilisées : RAG, LLM, validation de Chauvin, seuils, scénarios
  \item La structure complète des fichiers et dossiers du projet
  \item Les étapes de réalisation détaillées avec les commandes et résultats
  \item La veille technologique sur les outils et frameworks choisis
\end{itemize}

\section{Contexte du projet}

L'Agent Hypothèses s'inscrit dans la \textbf{Phase 2 — Stratégie \& Business Plan} de la plateforme. Son rôle est de conduire un dialogue pédagogique structuré avec l'entrepreneur pour collecter \textbf{22 hypothèses} organisées en 4 blocs, les valider selon les règles métier de Chauvin, et produire un JSON structuré transmis à l'Agent Finance pour la génération des tableaux prévisionnels.

\vspace{8pt}
\begin{center}
\begin{tabularx}{0.92\textwidth}{
  >{\bfseries\color{navyblue}}L{3.8cm}
  X
}
\toprule
\rowcolor{navyblue}\textcolor{white}{\textbf{Composante}} & \textcolor{white}{\textbf{Description}} \\
\midrule
\rowcolor{paleblue} Périmètre & Pré-création d'entreprise — Phase Hypothèses \\
LLM utilisé & Llama-3.3-70b-versatile via API Groq \\
\rowcolor{paleblue} Base vectorielle & ChromaDB avec 280 documents indexés \\
Langue & Français (Darija et Arabe en extension post-PFE) \\
\rowcolor{paleblue} Output & JSON structuré HypothesesOutput $\rightarrow$ Agent Finance \\
API & FastAPI REST sur port 8000 \\
\bottomrule
\end{tabularx}
\end{center}

\section{Architecture globale du workflow multi-agents}

\begin{formulebox}[title={\bfseries\color{navyblue} Workflow de la plateforme}]
\texttt{Angular $\rightarrow$ FastAPI $\rightarrow$ SLM Phi-3 Mini (classification intention)}\\
\texttt{$\rightarrow$ Agent Contrôleur n8n $\rightarrow$ POST /hypotheses/start}\\
\texttt{$\rightarrow$ Dialogue 22 questions (RAG Chauvin + Groq)}\\
\texttt{$\rightarrow$ Validation Chauvin $\rightarrow$ JSON HypothesesOutput}\\
\texttt{$\rightarrow$ Agent Finance (POST /plan) $\rightarrow$ Agent Business Plan $\rightarrow$ PDF final}
\end{formulebox}

\section{Structure des fichiers et démarche de réalisation}

La structure du projet respecte une architecture modulaire où chaque dossier correspond à une couche fonctionnelle distincte. Cette organisation garantit la lisibilité, la maintenabilité et la testabilité de chaque composant de manière indépendante.

\subsection{Arborescence du projet}

\begin{codebox}
agent\_hypotheses/\\
\quad{}+-- .env\\
\quad{}+-- main.py\\
\quad{}+-- requirements.txt\\
\quad{}+-- agent/\\
\quad{}|\quad{}+-- \_\_init\_\_.py\\
\quad{}|\quad{}+-- dialogue.py\\
\quad{}|\quad{}+-- groq\_client.py\\
\quad{}|\quad{}+-- validation.py\\
\quad{}|\quad{}+-- questions.py\\
\quad{}|\quad{}+-- scenarios.py\\
\quad{}|\quad{}\textbackslash{}-- output\_builder.py\\
\quad{}+-- rag/\\
\quad{}|\quad{}+-- \_\_init\_\_.py\\
\quad{}|\quad{}+-- vectorstore.py\\
\quad{}|\quad{}+-- embedder.py\\
\quad{}|\quad{}\textbackslash{}-- retriever.py\\
\quad{}+-- schemas/\\
\quad{}|\quad{}\textbackslash{}-- output\_schema.py\\
\quad{}+-- data/\\
\quad{}|\quad{}\textbackslash{}-- scenarios\_reference.json\\
\quad{}\textbackslash{}-- scripts/\\
\quad{}\quad{}\quad{}+-- test\_dialogue.py\\
\quad{}\quad{}\quad{}+-- test\_validation.py\\
\quad{}\quad{}\quad{}+-- test\_json\_final.py\\
\quad{}\quad{}\quad{}+-- test\_scenarios\_rag.py\\
\quad{}\quad{}\quad{}\textbackslash{}-- build\_scenarios\_rag.py
\end{codebox}

\subsection{Explication détaillée de chaque fichier — Démarche par étapes}

Les fichiers du projet sont présentés ici dans l'ordre chronologique de leur création et de leur rôle dans le cycle de réalisation. Chaque fichier correspond à une étape de la démarche Data Analytics.

\bigskip

%% --- Étape 1 ---
\begin{stepbox}{Étape 1 — \texttt{.env} \& \texttt{requirements.txt} : Configuration de l'environnement}
\textbf{.env} stocke les variables sensibles : \texttt{GROQ\_API\_KEY}, \texttt{CHROMA\_DB\_PATH} et les paramètres de connexion. Ce fichier est chargé au démarrage grâce à \texttt{load\_dotenv()} et n'est jamais versionné (exclu via \texttt{.gitignore}).

\medskip
\textbf{requirements.txt} liste toutes les dépendances nécessaires au projet : \texttt{fastapi}, \texttt{pydantic}, \texttt{python-dotenv}, \texttt{groq}, \texttt{chromadb}, \texttt{sentence-transformers}, \texttt{uvicorn}, \texttt{pypdf}, et \texttt{langchain}. Python 3.10 a été retenu pour sa compatibilité optimale avec l'ensemble de ces bibliothèques.

\begin{codebox}
\# Création de l'environnement virtuel Python 3.10\\
py -3.10 -m venv venv\\
.{\textbackslash}venv{\textbackslash}Scripts{\textbackslash}activate\\
pip install -r requirements.txt
\end{codebox}
\end{stepbox}

\bigskip

%% --- Étape 2 ---
\begin{stepbox}{Étape 2 — \texttt{schemas/output\_schema.py} : Le contrat JSON}
Ce fichier définit la structure exacte du JSON produit par l'agent à destination de l'Agent Finance. Il utilise \textbf{Pydantic} pour garantir la cohérence et le typage de chaque champ.

Les classes définies sont les suivantes :
\begin{itemize}
  \item \texttt{Bloc1Ventes} — hypothèses de ventes, clients, prix et croissance
  \item \texttt{Bloc2Achats} — type d'activité, coûts de fabrication, fournisseurs
  \item \texttt{Bloc3Charges} — charges fixes : loyer, salaires, utilités, marketing, investissements
  \item \texttt{Bloc4Encaissements} — nature clients et règle d'encaissement
  \item \texttt{ResultatValidation} — statut et messages des 3 règles de Chauvin
  \item \texttt{Scenario} — CA mensuel/annuel, charges, marge, résultat net, mois de rentabilité
  \item \texttt{TroisScenarios} — agrège les 3 projections (pessimiste, réaliste, optimiste)
  \item \texttt{HypothesesOutput} — classe principale qui agrège tous les blocs ci-dessus
\end{itemize}

Ce fichier est le premier livrable concret du projet : il constitue le \textbf{contrat formel} entre l'Agent Hypothèses et l'Agent Finance.
\end{stepbox}

\bigskip

%% --- Étape 3 ---
\begin{stepbox}{Étape 3 — \texttt{agent/questions.py} : Les 22 questions structurées}
Ce fichier contient le dictionnaire \texttt{QUESTIONS} définissant les 22 questions du dialogue, organisées en 4 blocs selon la méthode Chauvin. Chaque entrée précise le texte de la question, l'identifiant de la variable collectée, et le type de réponse attendu (choix, numérique, pourcentage, texte libre, etc.).

Ce fichier est entièrement statique : il n'appelle ni le LLM ni la base vectorielle. Il sert de référentiel structurel pour \texttt{dialogue.py}.
\end{stepbox}

\bigskip

%% --- Étape 4 ---
\begin{stepbox}{Étape 4 — \texttt{rag/vectorstore.py}, \texttt{rag/embedder.py}, \texttt{rag/retriever.py} : Pipeline RAG}
Ces trois fichiers constituent la couche \textbf{Retrieval-Augmented Generation} du projet.

\medskip
\textbf{vectorstore.py} expose la fonction \texttt{init\_chromadb(path)} qui ouvre — ou crée — la base ChromaDB persistante. La base contient 280 documents : 270 chunks extraits du livre Chauvin et 10 scénarios marocains de référence.

\medskip
\textbf{embedder.py} expose la fonction \texttt{get\_embedder()} qui instancie le modèle de représentation sémantique \texttt{all-MiniLM-L6-v2} (sentence-transformers / HuggingFace). Ce modèle léger de 80 Mo convertit chaque texte en un vecteur de 384 dimensions.

\medskip
\textbf{retriever.py} expose la fonction \texttt{search\_context(query, collection, embedder)} qui effectue la recherche par similarité cosinus dans ChromaDB et retourne les 3 passages les plus pertinents pour enrichir chaque prompt envoyé à Groq.

\begin{codebox}
\# Indexation du livre Chauvin\\
python scripts/build\_rag.py\\
\# → 78 pages extraites → 270 chunks vectorisés\\
\\
\# Indexation des scénarios marocains\\
python scripts/build\_scenarios\_rag.py\\
\# → 10 scénarios ajoutés → total 280 documents
\end{codebox}
\end{stepbox}

\bigskip

%% --- Étape 5 ---
\begin{stepbox}{Étape 5 — \texttt{agent/groq\_client.py} : Connexion au LLM}
Ce fichier expose la fonction \texttt{init\_groq()} qui instancie et configure le client Groq (modèle \texttt{llama-3.3-70b-versatile}, température et timeout). Il centralise l'accès à l'API LLM : tous les autres modules appellent ce client unique au lieu d'instancier leur propre connexion.

La connexion Groq répond en moyenne en \textbf{1~070~ms}, bien en-dessous du seuil de 2 secondes fixé dans les exigences du projet.
\end{stepbox}

\bigskip

%% --- Étape 6 ---
\begin{stepbox}{Étape 6 — \texttt{agent/validation.py} : Les règles de Chauvin}
Ce fichier contient toutes les règles de validation codées en Python pur, sans appel au LLM (validation instantanée). Il définit :

\begin{itemize}
  \item \textbf{SEUILS\_VARIABLES} — dictionnaire des bornes min/max de chaque hypothèse numérique (ex. prix de vente entre 1 et 1~000~000 MAD, taux de croissance entre 0 et 200~\%)
  \item \textbf{valider\_seuils()} — vérifie les bornes et renvoie le statut OK / AVERTISSEMENT / ERREUR pour chaque variable
  \item \textbf{valider\_seuil\_rentabilite()} — Règle~1 : compare les clients prévus au seuil minimum calculé
  \item \textbf{valider\_symetrie\_financement()} — Règle~2 : vérifie que les investissements ne dépassent pas le double des apports propres
  \item \textbf{valider\_coherence\_prix\_fabrication()} — Règle~3 : s'assure que le coût de fabrication reste sous 80~\% du prix de vente
  \item \textbf{valider\_toutes\_hypotheses()} — orchestre les trois règles et produit le dictionnaire de validation global transmis à l'output builder
\end{itemize}

\begin{codebox}
\# Test unitaire de la validation\\
python scripts/test\_validation.py\\
\# → Charges fixes : 10 742 MAD/mois\\
\# → Seuil clients minimum : 154 ventes/mois\\
\# → Clients prévus mois 1 : 30 → AVERTISSEMENT Chauvin
\end{codebox}
\end{stepbox}

\bigskip

%% --- Étape 7 ---
\begin{stepbox}{Étape 7 — \texttt{agent/scenarios.py} : Les 3 projections financières}
Ce fichier calcule les trois scénarios financiers à partir des hypothèses collectées. La fonction \texttt{calculer\_ca\_mensuel()} applique la formule de croissance composée :

$$\text{CA mois } N = \text{prix} \times \text{clients} \times \left(1 + \frac{\text{taux}}{100}\right)^N$$

La fonction \texttt{generer\_scenarios(hypotheses)} applique les coefficients 0{,}7 / 1{,}0 / 1{,}3 aux hypothèses de base et produit pour chaque scénario : le CA sur 12~mois, le CA annuel total, les charges annuelles, la marge brute, le résultat net et le mois de rentabilité (jusqu'au mois~24). Elle retourne un objet \texttt{TroisScenarios}.
\end{stepbox}

\bigskip

%% --- Étape 8 ---
\begin{stepbox}{Étape 8 — \texttt{agent/output\_builder.py} : Construction du JSON final}
Ce fichier contient la fonction \texttt{construire\_json\_final()} qui orchestre la production du livrable de l'agent :

\begin{enumerate}
  \item Mappe les valeurs collectées vers les sous-modèles Pydantic (Bloc1 à Bloc4)
  \item Détermine la règle d'encaissement selon la nature du client (B2B / B2C / mixte)
  \item Appelle \texttt{generer\_scenarios()} pour les 3 projections financières
  \item Calcule le message pédagogique final via \texttt{determiner\_message\_final()}
  \item Retourne une instance \texttt{HypothesesOutput} complète, validée par Pydantic
\end{enumerate}

\begin{codebox}
python scripts/test\_json\_final.py\\
\# → session\_id : cd342c82\\
\# → statut\_global : AVERTISSEMENTS\\
\# → 12 champs obligatoires présents : \checkmark{} JSON valide
\end{codebox}
\end{stepbox}

\bigskip

%% --- Étape 9 ---
\begin{stepbox}{Étape 9 — \texttt{agent/dialogue.py} : Le cœur conversationnel de l'agent}
Ce fichier contient la logique principale d'orchestration du dialogue. Il coordonne les 22 questions bloc par bloc, appelle le RAG pour enrichir chaque question avec du contexte Chauvin, détecte l'intention de chaque réponse via Groq (réponse directe / question / hors-sujet / vide), applique les guardrails de protection du flux, valide chaque valeur saisie via \texttt{valider\_seuils()}, puis déclenche la validation globale à la fin de la collecte avant d'appeler \texttt{construire\_json\_final()}.

Le dialogue gère un compteur de tentatives (maximum 3 par question) et assure que les questions posées par l'entrepreneur pendant le dialogue ne comptent pas comme des tentatives ratées.
\end{stepbox}

\bigskip

%% --- Étape 10 ---
\begin{stepbox}{Étape 10 — \texttt{main.py} : L'API REST FastAPI}
Point d'entrée de l'application. Ce fichier crée une instance FastAPI, initialise une fois le client Groq, la collection ChromaDB et l'embedder, puis expose 5 endpoints :

\begin{center}
\begin{tabularx}{0.88\textwidth}{
  >{\ttfamily\small\color{navyblue}}L{3.5cm}
  >{\centering\arraybackslash\small}C{1.5cm}
  X
  >{\centering\arraybackslash\small}C{2cm}
}
\toprule
\rowcolor{navyblue}\textcolor{white}{\textbf{Endpoint}} & \textcolor{white}{\textbf{Méthode}} & \textcolor{white}{\textbf{Description}} & \textcolor{white}{\textbf{Statut}} \\
\midrule
\rowcolor{paleblue}/health & GET & État du service + nb docs RAG & 200 OK \checkmark{} \\
/hypotheses/start & POST & Lance le dialogue complet & 200 OK \checkmark{} \\
\rowcolor{paleblue}/hypotheses/validate & POST & Valide sans dialogue & 200 OK \checkmark{} \\
/hypotheses/questions & GET & Liste des 22 questions & 200 OK \checkmark{} \\
\rowcolor{paleblue}/scenarios/\{id\} & GET & Historique de session & En attente (en attente) \\
\bottomrule
\end{tabularx}
\end{center}
\end{stepbox}

\bigskip

%% --- Étape 11 ---
\begin{stepbox}{Étape 11 — \texttt{scripts/} : Tests et scripts de validation}
Le dossier \texttt{scripts/} regroupe tous les outils de vérification indépendants du serveur :

\begin{itemize}
  \item \textbf{test\_dialogue.py} — Simule un dialogue complet (sans UI) : affiche le JSON final et les logs de validation pour un scénario de test prédéfini.
  \item \textbf{test\_validation.py} — Test unitaire : crée un jeu d'hypothèses, appelle \texttt{valider\_toutes\_hypotheses()} et affiche le résultat détaillé.
  \item \textbf{test\_json\_final.py} — Combine validation et construction du JSON final ; vérifie que chaque champ obligatoire est présent.
  \item \textbf{test\_scenarios\_rag.py} — Vérifie que la recherche RAG retourne des scénarios pertinents pour une requête donnée.
  \item \textbf{build\_scenarios\_rag.py} — Script d'initialisation (exécuté une seule fois) : parcourt \texttt{scenarios\_reference.json} et indexe les 10 scénarios dans ChromaDB.
\end{itemize}
\end{stepbox}

\bigskip

%% --- Étape 12 ---
\begin{stepbox}{Étape 12 — \texttt{data/scenarios\_reference.json} : Données de référence RAG}
Ce fichier JSON contient 10 scénarios synthétiques marocains 2024 (\texttt{synthetique\_maroc\_2024}), couvrant différents secteurs d'activité (épicerie, restauration, artisanat, services B2B, etc.). Ces scénarios servent à la fois à l'indexation RAG (permettant à l'agent de comparer les hypothèses de l'entrepreneur à des cas réels) et aux tests de non-régression du pipeline de recherche sémantique.
\end{stepbox}

\vspace{12pt}
\begin{resultbox}
\textbf{\color{darkgreen} Résumé de l'architecture en 4 couches :}\\
\textbf{API} (\texttt{main.py}) : point d'entrée HTTP, orchestre les appels aux modules.\\
\textbf{Agent} (\texttt{agent/}) : dialogue, validation métier et construction du modèle de sortie.\\
\textbf{RAG} (\texttt{rag/}) : recherche de contexte dans une base vectorielle pour enrichir les prompts.\\
\textbf{Schémas} (\texttt{schemas/}) : description formelle du format JSON attendu.
\end{resultbox}

%% ============================================================
%%  CHAPITRE 2 — NOTIONS FONDAMENTALES & LOGIQUE MÉTIER
%% ============================================================
\chapter{Notions Fondamentales \& Logique Métier}

\section{LLM — Large Language Model}

Un LLM est un modèle de langage de grande taille entraîné sur des milliards de tokens textuels. Il est capable de comprendre le langage naturel et de générer du texte cohérent en réponse à une instruction (prompt). Dans ce projet, le LLM joue le rôle du cerveau conversationnel de l'agent.

\begin{formulebox}[title={\bfseries\color{navyblue} Rôle du LLM dans l'Agent Hypothèses}]
Le LLM \textbf{Llama-3.3-70b-versatile}, accessible via l'API Groq, remplit cinq fonctions :
\begin{itemize}
  \item Génère les questions pédagogiques en français naturel
  \item Explique chaque hypothèse à l'entrepreneur
  \item Valide les réponses contextuelles
  \item Répond aux questions posées pendant le dialogue
  \item Détecte l'intention derrière chaque réponse (réponse directe, question, hors-sujet, vide)
\end{itemize}
\end{formulebox}

\section{RAG — Retrieval-Augmented Generation}

Le RAG est une architecture qui combine la recherche sémantique dans une base de données vectorielle avec la génération de texte par un LLM. Au lieu que le LLM réponde uniquement depuis sa mémoire interne (risque d'hallucination), il récupère d'abord les passages les plus pertinents depuis une base de connaissances externe, puis génère une réponse ancrée dans des sources vérifiées.

Le pipeline RAG de l'agent fonctionne en trois étapes :

\begin{enumerate}
  \item La question ou le contexte est converti en vecteur numérique (\textit{embedding}) par le modèle \texttt{all-MiniLM-L6-v2}
  \item ChromaDB compare ce vecteur avec les 280 documents indexés par similarité cosinus et retourne les 3 passages les plus proches
  \item Ces passages sont injectés dans le prompt envoyé à Groq, qui génère une réponse enrichie par les sources Chauvin et les scénarios marocains
\end{enumerate}

\section{Embeddings et Similarité Cosinus}

Un \textit{embedding} est la représentation numérique d'un texte sous forme d'un vecteur de 384 dimensions. Deux textes qui parlent du même sujet ont des vecteurs proches dans l'espace vectoriel, même s'ils utilisent des mots différents. La similarité cosinus mesure l'angle entre deux vecteurs — plus l'angle est petit, plus les textes sont sémantiquement similaires.

\begin{formulebox}[title={\bfseries\color{navyblue} Modèle d'embedding utilisé}]
\textbf{all-MiniLM-L6-v2} (sentence-transformers / HuggingFace) — modèle léger de 80~Mo, optimisé pour la recherche sémantique. Transforme chaque chunk de texte en un vecteur de 384 nombres. Produit \textbf{270 vecteurs} depuis le livre Chauvin et \textbf{10 vecteurs} depuis les scénarios marocains de référence.
\end{formulebox}

\section{ChromaDB — Base Vectorielle}

ChromaDB est une base de données vectorielle open-source qui stocke les embeddings et permet leur recherche par similarité sémantique.

\begin{center}
\begin{tabularx}{0.88\textwidth}{X X}
\toprule
\rowcolor{navyblue}\textcolor{white}{\textbf{PostgreSQL}} & \textcolor{white}{\textbf{ChromaDB}} \\
\midrule
\rowcolor{paleblue} Base relationnelle & Base vectorielle \\
Stocke lignes et colonnes & Stocke textes et vecteurs \\
\rowcolor{paleblue} Requête SQL & Requête par similarité sémantique \\
Résultat exact & Résultat par pertinence \\
\rowcolor{paleblue} Sessions utilisateurs, logs & Livre Chauvin, scénarios marocains \\
\bottomrule
\end{tabularx}
\end{center}

\section{Les Hypothèses Financières — Méthode Chauvin}

Une hypothèse financière est une supposition raisonnée et quantifiée qui sert de base de calcul pour les projections financières d'un Business Plan. Comme le souligne Pascal Chauvin, l'imprécision n'a pas de place dans le domaine financier : chaque hypothèse doit être réaliste, modeste et défendable.

Le module s'appuie sur les 7 questions structurantes de Chauvin (\textit{Qui~? Quoi~? Comment~? Pourquoi~? Quand~? Où~? Combien~?}), dont la question \textbf{COMBIEN~?} est directement traduite en 4 blocs d'hypothèses à collecter.

\section{Les 4 Blocs d'Hypothèses}

\begin{center}
\begin{tabularx}{0.95\textwidth}{
  >{\centering\arraybackslash\bfseries\color{navyblue}}C{1.2cm}
  >{\centering\arraybackslash}C{2.2cm}
  X
  >{\centering\arraybackslash}C{2cm}
}
\toprule
\rowcolor{navyblue}\textcolor{white}{\textbf{Bloc}} & \textcolor{white}{\textbf{Question Chauvin}} & \textcolor{white}{\textbf{Contenu}} & \textcolor{white}{\textbf{Hypothèses}} \\
\midrule
\rowcolor{paleblue} 1 & Quoi ? & Ventes, clients, prix, croissance, fidélisation & H1 à H7 \\
2 & Comment ? & Type activité, fabrication, infrastructure, fournisseurs & H8 à H12 \\
\rowcolor{paleblue} 3 & Combien ? & Charges fixes, loyer, salaires, marketing, investissements & H13 à H21 \\
4 & Quand ? & Nature clients, délais encaissement & H22 \\
\bottomrule
\end{tabularx}
\end{center}

\section{Règles de Validation de Chauvin}

\subsection{Règle 1 — Seuil de Rentabilité}

Le seuil de rentabilité est le niveau de ventes minimum pour que l'entreprise couvre exactement ses charges fixes sans réaliser de bénéfice ni de perte.

\begin{formulebox}[title={\bfseries\color{navyblue} Formule du seuil de rentabilité}]
\[
\text{Charges fixes} = H_{13} + H_{14} + H_{15} + H_{16} + H_{17} + \frac{H_{18}}{12}
\]
\[
\text{Marge unitaire} = H_2 \,(\text{prix de vente}) - H_9 \,(\text{coût de fabrication})
\]
\[
\text{Seuil clients minimum} = \frac{\text{Charges fixes}}{\text{Marge unitaire}}
\]
\textbf{Si clients prévus $<$ seuil} $\Rightarrow$ \textcolor{darkgreen}{\textbf{AVERTISSEMENT pédagogique}}
\end{formulebox}

\subsection{Règle 2 — Symétrie du Financement}

Selon Chauvin, si les besoins de financement dépassent le double des apports personnels de l'entrepreneur, la structure financière est déséquilibrée.

\begin{formulebox}[title={\bfseries\color{navyblue} Formule de la symétrie}]
\[
\text{Apports propres} = H_{19} - H_{21}
\]
\textbf{Si } $H_{19} > 2 \times \text{apports propres}$ $\Rightarrow$ \textcolor{darkgreen}{\textbf{AVERTISSEMENT Chauvin}}
\end{formulebox}

\subsection{Règle 3 — Cohérence Prix / Coût de Fabrication}

Si le coût de fabrication unitaire représente plus de 80~\% du prix de vente, la marge brute est insuffisante pour couvrir les charges fixes. L'agent alerte l'entrepreneur et suggère d'augmenter le prix ou de réduire les coûts.

\section{Les 3 Scénarios Financiers}

Conformément à la méthode Chauvin, l'agent génère automatiquement trois versions du prévisionnel financier en appliquant des coefficients multiplicateurs aux hypothèses de base.

\begin{center}
\begin{tabularx}{0.95\textwidth}{
  >{\bfseries\color{navyblue}}L{2.5cm}
  >{\centering\arraybackslash}C{2cm}
  X
  L{4cm}
}
\toprule
\rowcolor{navyblue}\textcolor{white}{\textbf{Scénario}} & \textcolor{white}{\textbf{Coefficient}} & \textcolor{white}{\textbf{Logique}} & \textcolor{white}{\textbf{Usage}} \\
\midrule
\rowcolor{paleblue} Pessimiste & $\times\, 0{,}7$ & 30~\% moins de clients & Vérifier la survie \\
Réaliste & $\times\, 1{,}0$ & Hypothèses telles que saisies & Prévisions de référence \\
\rowcolor{paleblue} Optimiste & $\times\, 1{,}3$ & 30~\% plus de clients & Capacité à croître \\
\bottomrule
\end{tabularx}
\end{center}

Pour chaque scénario, l'agent calcule : le CA mensuel sur 12~mois (croissance composée), le CA annuel total, les charges annuelles, la marge brute, le résultat net annuel et le mois de rentabilité.

\section{Guardrails — Gestion des Intentions}

Les guardrails sont des mécanismes de protection qui contrôlent le flux du dialogue et empêchent l'agent de se bloquer face à des réponses inattendues. L'agent détecte 4 types d'intention :

\begin{center}
\begin{tabularx}{0.95\textwidth}{
  >{\bfseries\color{navyblue}}L{2.5cm}
  X
  L{4.5cm}
}
\toprule
\rowcolor{navyblue}\textcolor{white}{\textbf{Intention}} & \textcolor{white}{\textbf{Description}} & \textcolor{white}{\textbf{Action de l'agent}} \\
\midrule
\rowcolor{paleblue} Réponse & L'entrepreneur répond directement & Valider et enregistrer \\
Question & L'entrepreneur pose une question & Répondre via RAG puis reposer la question \\
\rowcolor{paleblue} Hors sujet & Réponse sans rapport & Signaler et reposer (max 3 tentatives) \\
Vide & Entrée vide ou espaces & Rappeler poliment et reposer \\
\bottomrule
\end{tabularx}
\end{center}

\section{Délégation vers les Autres Agents}

Quand l'entrepreneur pose une question hors du périmètre du livre Chauvin, l'agent détecte la catégorie et indique quel agent spécialisé devrait répondre.

\begin{center}
\begin{tabularx}{0.95\textwidth}{
  >{\ttfamily\small\color{navyblue}}L{2.2cm}
  X
  >{\bfseries}L{3.2cm}
}
\toprule
\rowcolor{navyblue}\textcolor{white}{\textbf{Catégorie}} & \textcolor{white}{\textbf{Exemple}} & \textcolor{white}{\textbf{Agent délégué}} \\
\midrule
\rowcolor{paleblue} chauvin & C'est quoi le seuil de rentabilité~? & Agent Hypothèses \\
general & C'est quoi une SARL~? & Agent Général \\
\rowcolor{paleblue} marche & Mon secteur est-il saturé~? & Agent Marché \\
finance & Comment calculer mon BFR~? & Agent Finance \\
\rowcolor{paleblue} hors\_sujet & Quel temps fait-il à Tanger~? & Refus poli \\
\bottomrule
\end{tabularx}
\end{center}

\section{Résultats obtenus — Validation du système}

\subsection{Métriques de performance}

\begin{center}
\begin{tabularx}{0.95\textwidth}{X >{\centering\arraybackslash}C{3cm} >{\centering\arraybackslash}C{3cm}}
\toprule
\rowcolor{navyblue}\textcolor{white}{\textbf{Indicateur}} & \textcolor{white}{\textbf{Valeur obtenue}} & \textcolor{white}{\textbf{Statut}} \\
\midrule
\rowcolor{paleblue} Connexion Groq & 1~070~ms & \checkmark{} $<$ 2~s \\
Documents RAG & 280 indexés & \checkmark{} \\
\rowcolor{paleblue} Questions couvertes & 22 / 22 & \checkmark{} \\
Endpoints API & 5 / 5 & \checkmark{} \\
\rowcolor{paleblue} Scénarios générés & 3 automatiques & \checkmark{} \\
JSON final & 12 champs validés & \checkmark{} \\
\bottomrule
\end{tabularx}
\end{center}

\subsection{Exemple de résultat — Scénario épicerie Tanger}

Hypothèses saisies : H2 = 150~MAD (prix), H4 = 30~clients, H9 = 80~MAD (coût fabrication), H13 = 3~000~MAD (loyer), H14 = 6~000~MAD (salaires).

\begin{center}
\begin{tabularx}{0.95\textwidth}{X >{\centering\arraybackslash}C{2.8cm} >{\centering\arraybackslash}C{2.8cm} >{\centering\arraybackslash}C{2.8cm}}
\toprule
\rowcolor{navyblue}\textcolor{white}{\textbf{Indicateur}} & \textcolor{white}{\textbf{Pessimiste}} & \textcolor{white}{\textbf{Réaliste}} & \textcolor{white}{\textbf{Optimiste}} \\
\midrule
\rowcolor{paleblue} Clients mois~1 & 21 & 30 & 39 \\
CA mois~1 & 3~150~MAD & 4~500~MAD & 5~850~MAD \\
\rowcolor{paleblue} CA annuel & 67~360~MAD & 96~229~MAD & 125~098~MAD \\
Charges annuelles & \multicolumn{3}{c}{143~900~MAD} \\
\rowcolor{paleblue} Résultat net an~1 & $-$112~465~MAD & $-$98~993~MAD & $-$85~521~MAD \\
Mois rentabilité & Mois~22 & Mois~19 & Mois~16 \\
\bottomrule
\end{tabularx}
\end{center}

\begin{resultbox}
\textbf{\color{darkgreen} Message pédagogique généré~:} Votre projet nécessite 19~mois avant rentabilité dans le scénario réaliste. Assurez-vous d'avoir une trésorerie suffisante pour couvrir les 19~premiers mois d'activité.
\end{resultbox}

%% ============================================================
%%  CHAPITRE 3 — VEILLE TECHNOLOGIQUE & STACK TECHNIQUE
%% ============================================================
\chapter{Veille Technologique \& Stack Technique de l'Agent}

\section{Vue d'ensemble}

Cette section présente les technologies choisies pour construire l'Agent Hypothèses, avec pour chaque outil une justification du choix par rapport aux alternatives disponibles sur le marché.

\section{Groq API — LLM haute performance}

Groq est un service d'inférence LLM ultra-rapide basé sur des puces matérielles spécialisées appelées LPU (\textit{Language Processing Unit}). Contrairement aux services GPU classiques, les LPU sont optimisés pour l'inférence séquentielle des tokens, réduisant la latence à 1 à 2 secondes par requête.

\begin{center}
\begin{tabularx}{0.95\textwidth}{X >{\centering\arraybackslash}C{2.8cm} >{\centering\arraybackslash}C{2.8cm} >{\centering\arraybackslash}C{2.8cm}}
\toprule
\rowcolor{navyblue}\textcolor{white}{\textbf{Critère}} & \textcolor{white}{\textbf{Groq API}} & \textcolor{white}{\textbf{OpenAI API}} & \textcolor{white}{\textbf{Ollama (local)}} \\
\midrule
\rowcolor{paleblue} Vitesse & Ultra-rapide (LPU) & Rapide (GPU) & Dépend du matériel \\
Coût & Plan gratuit & Payant & Gratuit \\
\rowcolor{paleblue} Modèle 70B & llama-3.3-70b & GPT-4o & llama3:70b \\
Confidentialité & API externe & API externe & 100~\% local \\
\rowcolor{lightgreen} \textbf{Choix PFE} & \textbf{\checkmark{} Sélectionné} & Trop coûteux & GPU insuffisant \\
\bottomrule
\end{tabularx}
\end{center}

\section{ChromaDB — Base Vectorielle}

ChromaDB est la base vectorielle open-source la plus adoptée pour les projets RAG en Python, avec une intégration native LangChain. Elle supporte le stockage persistant, la recherche par similarité cosinus et le filtrage par métadonnées.

\begin{formulebox}[title={\bfseries\color{navyblue} Comparaison des bases vectorielles}]
\begin{itemize}
  \item \textbf{ChromaDB (retenu)} : open-source, Python natif, stockage local, gratuit, intégration LangChain directe
  \item \textbf{Pinecone} : cloud uniquement, performances supérieures mais coûteux
  \item \textbf{Weaviate} : plus complexe, adapté aux grands volumes
  \item \textbf{FAISS (Meta)} : très performant mais sans API native de persistance
\end{itemize}
ChromaDB est optimal pour un PFE avec contraintes de budget.
\end{formulebox}

\section{Sentence-Transformers — Modèle d'embedding}

La bibliothèque sentence-transformers de HuggingFace fournit des modèles pré-entraînés pour la conversion de texte en vecteurs. Le modèle \texttt{all-MiniLM-L6-v2} a été sélectionné pour son équilibre entre taille (80~Mo), vitesse et qualité de représentation sémantique.

\section{FastAPI — Framework API REST}

FastAPI est le framework Python le plus performant pour créer des APIs REST. Il génère automatiquement la documentation Swagger (accessible sur \texttt{/docs}), valide les données entrantes via Pydantic, et supporte les opérations asynchrones pour une meilleure performance.

\section{Pydantic — Validation et Sérialisation}

Pydantic est la bibliothèque Python standard pour la validation de données et la définition de schémas. Elle garantit que le JSON produit par l'agent respecte exactement la structure attendue par l'Agent Finance, avec typage fort et messages d'erreur clairs.

\section{LangChain — Orchestration RAG}

LangChain est le framework d'orchestration des pipelines LLM le plus utilisé dans l'industrie. Il simplifie la connexion entre le LLM (Groq), la base vectorielle (ChromaDB) et le module d'embedding (sentence-transformers), en fournissant des abstractions prêtes à l'emploi pour les chaînes de raisonnement RAG.

\section{n8n — Orchestration Multi-Agents}

n8n est la plateforme d'orchestration de workflows open-source qui joue le rôle de chef d'orchestre de la plateforme complète. Il reçoit les requêtes depuis Angular (via FastAPI), classe l'intention avec le SLM Phi-3~Mini, route vers l'agent spécialiste approprié, passe le JSON de l'Agent Hypothèses à l'Agent Finance, et assemble la réponse finale. L'intégration de l'Agent Hypothèses dans n8n se fait via l'endpoint \texttt{POST /hypotheses/start}.

\section{Stack Technologique Complète}

\begin{center}
\begin{tabularx}{0.95\textwidth}{
  >{\bfseries\color{navyblue}}L{3cm}
  L{3.5cm}
  X
}
\toprule
\rowcolor{navyblue}\textcolor{white}{\textbf{Catégorie}} & \textcolor{white}{\textbf{Technologie}} & \textcolor{white}{\textbf{Version / Rôle}} \\
\midrule
\rowcolor{paleblue} Langage & Python & 3.10.11 — compatibilité optimale \\
LLM & Groq API & llama-3.3-70b-versatile \\
\rowcolor{paleblue} RAG & LangChain & Orchestration pipeline RAG \\
Base vectorielle & ChromaDB & PersistentClient — 280 docs \\
\rowcolor{paleblue} Embeddings & sentence-transformers & all-MiniLM-L6-v2 — 384 dims \\
Extraction PDF & pypdf & 78 pages → 270 chunks \\
\rowcolor{paleblue} API REST & FastAPI + uvicorn & Port 8000 \\
Validation & Pydantic & HypothesesOutput schema \\
\rowcolor{paleblue} Orchestration & n8n & Connexion finale multi-agents \\
Environnement & venv Python 3.10 & Isolation des dépendances \\
\bottomrule
\end{tabularx}
\end{center}

\section{Conclusion et Perspectives}

L'Agent Hypothèses constitue le pont pédagogique entre l'entrepreneur et les outils analytiques de la plateforme. En collectant les 22~hypothèses à travers un dialogue conversationnel ancré dans la méthode Chauvin, en les validant avec des règles codées en dur et en produisant un JSON structuré prêt pour l'Agent Finance, cet agent remplit exactement le rôle défini dans le Cahier des Charges : accompagner l'entrepreneur de l'idée aux hypothèses validées, de manière pédagogique, fiable et multilingue.

\vspace{10pt}
\begin{center}
\begin{tabularx}{0.88\textwidth}{X X}
\toprule
\rowcolor{navyblue}\textcolor{white}{\textbf{Ce qui est accompli \checkmark{}}} & \textcolor{white}{\textbf{Prochaines étapes (en attente)}} \\
\midrule
\rowcolor{paleblue} 22 hypothèses en 4 blocs & Connexion PostgreSQL sessions \\
RAG 280 documents Chauvin & Extension Darija et Arabe \\
\rowcolor{paleblue} Validation 3 règles Chauvin & Intégration n8n workflow \\
3 scénarios automatiques & Golden Set 50 scénarios \\
\rowcolor{paleblue} API FastAPI 5 endpoints & Tests avec profils réels \\
JSON contrat Agent Finance & Module OCR formulaires \\
\bottomrule
\end{tabularx}
\end{center}

\end{document}
